Penelitian ini mengkaji potensi ketidaksesuaian antara formulasi prediksi satu-langkah dalam \textit{opponent modelling} dan kebutuhan perencanaan strategis multi-langkah dalam permainan berulang seperti \textit{Iterated Prisoner's Dilemma}. Melalui implementasi modul prediksi yang difokuskan pada \textit{k-step forecasting}, penelitian ini mengevaluasi bagaimana horizon prediksi memengaruhi akurasi estimasi distribusi aksi lawan serta dinamika kesalahan prediksi antar-langkah.

Hasil eksperimen menunjukkan bahwa perluasan horizon prediksi memberikan representasi yang lebih informatif terhadap pola strategi lawan yang bergantung pada beberapa langkah interaksi, meskipun tetap terdapat trade-off dalam bentuk peningkatan kompleksitas dan potensi propagasi kesalahan. Temuan ini mengindikasikan bahwa pemilihan objektif pembelajaran dalam \textit{opponent modelling} perlu mempertimbangkan keselarasan dengan horizon perencanaan strategis yang digunakan agen.

Secara keseluruhan, penelitian ini memberikan dasar eksperimental untuk memahami peran horizon prediksi dalam interaksi strategis berulang, tanpa melakukan asumsi terhadap struktur internal atau parametrik dari algoritma lawan.